\documentclass[11pt,a4paper]{article}
\usepackage[utf8]{inputenc}
\usepackage[T1]{fontenc}
\usepackage[portuguese]{babel}
\usepackage{xcolor}
\usepackage{hyperref}
\usepackage{geometry}
\usepackage{tcolorbox}

\geometry{margin=2.5cm}

\definecolor{primary}{RGB}{41, 128, 185}
\definecolor{secondary}{RGB}{44, 62, 80}
\definecolor{codebg}{RGB}{240, 240, 240}

\hypersetup{colorlinks=true, linkcolor=primary, urlcolor=primary}

\title{\color{secondary}\Huge\textbf{Exercício 5.6: Testando Serverless}}
\author{\color{primary}DevOps com Kubernetes -- 2026}
\date{}

\begin{document}
\maketitle

\section*{\color{primary}Instruções do Exercício}
\begin{tcolorbox}[colback=blue!5!white,colframe=blue!75!black,title=Exercício 5.6: Testando Serverless]
Experimente os conceitos de \textit{Serverless computing} criando sua própria infraestrutura em cima do Kubernetes.

- Instale o componente \textbf{Knative Serving} no seu cluster k3d (lembre-se de criar o cluster desabilitando o roteador Traefik nativo para não haver conflito de portas).

- Siga o guia oficial de instalação via YAML, escolhendo o Kourier como camada de rede e o Magic DNS para resolução local.

- Se houver falhas e \textit{CrashLoopBackOffs}, inspecione os logs para corrigi-los.

- Por fim, execute e teste os exemplos de Deploy, Autoscaling (escalonamento a zero) e Traffic splitting fornecidos na documentação do Knative.
\end{tcolorbox}

\end{document}
